% Problem-section file following ../_template.tex.
\section{Quantum capacity of a qubit Pauli channel}
\label{sec:problem-1}

\paragraph{Problem.}
What is the quantum capacity $\mathcal{Q}(\Lambda_{\mathbf p})$ of the qubit
Pauli channel defined by
\begin{equation}
  \Lambda_{\mathbf p}(\rho)
  =p_I\rho+p_XX\rho X+p_YY\rho Y+p_ZZ\rho Z,
  \qquad p_I+p_X+p_Y+p_Z=1?
  \label{eq:p1-pauli-channel}
\end{equation}
Equation~\eqref{eq:p1-pauli-channel} fixes the channel and the error-probability
convention used throughout this section.

\subsection*{Status}

Unsolved

\subsection*{Source}

The exact-capacity question is implicit in the hashing lower bound of Bennett
et al. and the strict degenerate-code improvement of DiVincenzo, Shor, and
Smolin, which together leave the general Pauli-channel capacity undetermined
\sourcecite{ref:p1-bennett}{BDSW96},
\sourcecite{ref:p1-divincenzo}{DSS98}.

\subsection*{Progress}

\begin{itemize}
  \item The hashing construction gives the achievable lower bound
  \begin{equation}
    \mathcal{Q}(\Lambda_{\mathbf p})\ge \max\{0,1-H(\mathbf p)\},
    \qquad
    H(\mathbf p):=-\sum_{i\in\{I,X,Y,Z\}}p_i\log_2p_i,
    \qquad 0\log_2 0:=0.
    \label{eq:p1-hashing-bound}
  \end{equation}
  The quantity $1-H(\mathbf p)$ in Eq.~\eqref{eq:p1-hashing-bound} equals
  $I_c(I/2,\Lambda_{\mathbf p})$, the coherent information at the maximally
  mixed input (often called the symmetric coherent information), rather than
  the optimized one-shot coherent information in general
  \sourcecite{ref:p1-bennett}{BDSW96}, \sourcecite{ref:p1-devetak}{Dev05}.

  \item The coherent-information rate in Eq.~\eqref{eq:p1-hashing-bound} is
  achieved by efficient quantum polar codes for qubit Pauli channels.  The
  first construction uses preshared entanglement (with zero consumption for
  sufficiently low noise), while a later construction removes the need for
  preshared entanglement \sourcecite{ref:p1-renes-2012}{RDR12},
  \sourcecite{ref:p1-renes-2015}{RSDR15}.

  \item In the depolarizing case, write $p_I=f$ and
  $p_X=p_Y=p_Z=(1-f)/3$, and let $r:=(1-f)/3$ be the per-Pauli error
  probability.  Krohn-Grimberghe constructed an explicit rank-two,
  permutation-symmetric input across $45$ channel uses and certified in exact
  rational arithmetic that its coherent information is positive at
  $r=16239/250000=0.064956$.  Monotonicity under channel post-processing
  therefore gives $\mathcal Q(\Lambda_{\mathbf p})>0$ whenever
  $f\ge201283/250000=0.805132$.  In particular, throughout
  $0.805132\le f<0.81071$ the capacity is positive although the hashing
  expression in Eq.~\eqref{eq:p1-hashing-bound} is negative.  Table~1 of the
  paper lists the preceding best printed positivity point as $r=0.064657$,
  obtained by Agarwal et al. through symmetry-reduced coherent-information
  optimization \sourcecite{ref:p1-agarwal-et-al}{AKL+26}.  The new point is
  machine-checkable but is not claimed to be the exact capacity threshold
  \sourcecite{ref:p1-krohn-grimberghe}{KG26}.

  \item For the qubit depolarizing subfamily, numerical optimization over
  channel extensions parameterized by the Stiefel manifold produces upper
  bounds on $\mathcal Q(\Lambda_{\mathbf p})$ that strictly improve the
  previously best flagged-extension bounds over the high-noise interval
  studied.  The optimized bounds do not meet the known achievable rates, so
  they narrow the capacity gap without determining the capacity.  The same
  work also proves that its amortized channel coherent information equals the
  ordinary one-shot channel coherent information for every channel, ruling
  out that amortization as a route to a stronger capacity lower bound
  \sourcecite{ref:p1-zhu-mao-fang-wang}{ZMFW25}.

  \item Kondra et al. obtained an efficiently computable, all-code
  exponential strong-converse bound for every full-support qubit Pauli
  channel.  For the depolarizing convention above, with total error
  $\delta:=1-f$, their bound reduces to
  $\mathcal Q(\Lambda_{\mathbf p})\le\max\{1-4\delta,0\}$; it is stricter
  than the regularized channel Rains bound for $\delta\gtrsim0.15642$ and
  vanishes at the antidegradability threshold $\delta=1/4$.  This sharpens
  the high-noise converse but does not determine the capacity below that
  threshold \sourcecite{ref:p1-kondra-et-al}{KBK+26}.
\end{itemize}

\subsection*{References}

\begin{enumerate}[label={},leftmargin=4.8em,labelsep=0.5em]
  \scriptsize
  \item[\textup{[BDSW96]}]\label{ref:p1-bennett}
  C. H. Bennett, D. P. DiVincenzo, J. A. Smolin, and W. K. Wootters,
  ``Mixed-State Entanglement and Quantum Error Correction,''
  \emph{Physical Review A} \textbf{54}, 3824--3851 (1996).
  \href{https://doi.org/10.1103/PhysRevA.54.3824}{doi:10.1103/PhysRevA.54.3824};
  \href{https://arxiv.org/abs/quant-ph/9604024}{arXiv:quant-ph/9604024}.

  \item[\textup{[Dev05]}]\label{ref:p1-devetak}
  I. Devetak, ``The Private Classical Capacity and Quantum Capacity of a
  Quantum Channel,'' \emph{IEEE Transactions on Information Theory}
  \textbf{51}, 44--55 (2005).
  \href{https://doi.org/10.1109/TIT.2004.839515}{doi:10.1109/TIT.2004.839515};
  \href{https://arxiv.org/abs/quant-ph/0304127}{arXiv:quant-ph/0304127}.

  \item[\textup{[RDR12]}]\label{ref:p1-renes-2012}
  J. M. Renes, F. Dupuis, and R. Renner,
  ``Efficient Polar Coding of Quantum Information,''
  \emph{Physical Review Letters} \textbf{109}, 050504 (2012).
  \newline
  \href{https://doi.org/10.1103/PhysRevLett.109.050504}{doi:10.1103/PhysRevLett.109.050504};
  \href{https://arxiv.org/abs/1109.3195}{arXiv:1109.3195}.

  \item[\textup{[RSDR15]}]\label{ref:p1-renes-2015}
  J. M. Renes, D. Sutter, F. Dupuis, and R. Renner, ``Efficient Quantum Polar
  Codes Requiring No Preshared Entanglement,'' \emph{IEEE Transactions on
  Information Theory} \textbf{61}, 6395--6414 (2015).
  \href{https://doi.org/10.1109/TIT.2015.2468084}{doi:10.1109/TIT.2015.2468084};
  \href{https://arxiv.org/abs/1307.1136}{arXiv:1307.1136}.

  \item[\textup{[DSS98]}]\label{ref:p1-divincenzo}
  D. P. DiVincenzo, P. W. Shor, and J. A. Smolin, ``Quantum-Channel Capacity
  of Very Noisy Channels,'' \emph{Physical Review A} \textbf{57}, 830--839
  (1998). \href{https://doi.org/10.1103/PhysRevA.57.830}{doi:10.1103/PhysRevA.57.830};
  \href{https://arxiv.org/abs/quant-ph/9706061}{arXiv:quant-ph/9706061}.

  \item[\textup{[AKL+26]}]\label{ref:p1-agarwal-et-al}
  A. Agarwal, A. R. Kalra, S. Lee, D. Leung, L. Schaeffer, P. Sinha, and
  G. Smith, ``Enhanced Quantum Capacity Thresholds from Symmetry,'' arXiv
  preprint (2026).
  \href{https://arxiv.org/abs/2605.09138}{arXiv:2605.09138}.

  \item[\textup{[KG26]}]\label{ref:p1-krohn-grimberghe}
  A. Krohn-Grimberghe,
  ``A Certified Lower Bound on the Quantum-Capacity Threshold of the
  Depolarizing Channel,'' arXiv preprint (2026).
  \href{https://arxiv.org/abs/2608.15870}{arXiv:2608.15870v2}.

  \item[\textup{[ZMFW25]}]\label{ref:p1-zhu-mao-fang-wang}
  C. Zhu, H. Mao, K. Fang, and X. Wang,
  ``Geometric Optimization for Quantum Communication,'' arXiv preprint
  (2025).
  \href{https://arxiv.org/abs/2509.15106v1}{arXiv:2509.15106}.

  \item[\textup{[KBK+26]}]\label{ref:p1-kondra-et-al}
  T. V. Kondra, R. Brinster, H. Kampermann, D. Bru{\ss}, and N. Wyderka,
  ``Sharp Quantum Capacity Thresholds: Exponential Strong Converses for
  Degradable and Antidegradable Channels,'' arXiv preprint (2026).
  \href{https://arxiv.org/abs/2608.01308}{arXiv:2608.01308}.
\end{enumerate}

\subsection*{Comment}

The exact quantum capacity remains unknown for a general qubit Pauli channel;
the hashing rate is an achievable lower bound and need not be the capacity.

\subsection*{Tag}

Pauli channels; Quantum error correction; Quantum capacity; Quantum Shannon
theory; Coherent information; Qubit systems

\subsection*{ID}

\texttt{op\_43d4aca67bd52554}
